\section{Validation Boundaries}
\label{sec:validation-boundaries}

An artificial society is not validated by the realism of a transcript, by one
aggregate match, or by successful execution.  GAWorld v2 separates five claim
levels: individual consistency, macro correspondence, emergent dynamics,
counterfactual response, and reproducibility.  Evidence at one level does not
compensate for absence at another.  This layered view follows established
agent-based-modeling practice in distinguishing implementation verification,
empirical validation, and the intended domain of use
\citep{grimm2006odd,windrum2007empirical,fagiolo2007critical}, and it now
extends the v1 system paper's framework with a sixth column for the
participant-evaluation claim level and an explicit evidence-class column
for the per-claim ledger.

\begin{table}[t]
\centering
\small
\caption{Claim--evidence matrix for the audited v2 snapshot.  ``Current
support'' states the strongest defensible claim, not an aspirational goal.
The ``Evidence class'' column is the class of the strongest available
artifact, not a statement that the strongest artifact is also a population
claim.}
\label{tab:claim-evidence}
\begin{tabular}{@{}p{0.16\linewidth}p{0.25\linewidth}p{0.21\linewidth}p{0.21\linewidth}p{0.13\linewidth}@{}}
\toprule
Claim level & Evidence required & Current support & Not established & Class \\
\midrule
Individual consistency & Repeated independent elicitations; longitudinal and perturbation tests; external or held-out criteria & Persistent profiles, memories, diaries, interviews, and one completed two-phase arm are inspectable & Human fidelity, stable identity across models/seeds, or a memory-treatment effect & \textsf{REAL} / \textsf{INCOMPLETE} \\
Macro correspondence & Pre-specified external targets; calibrated population; uncertainty and out-of-sample tests & State/economy pipelines can produce aggregates; benchmark fixtures test analysis code & Population representativeness, calibrated macro fit, forecasting, or real-city prediction & \textsf{REAL} / \textsf{FIXTURE} \\
Emergent dynamics & Multiple seeds and scales; mechanism ablations; network/temporal diagnostics; robust qualitative patterns & Social and network mechanisms emit inspectable state; failures are archived & Polarization, contagion, homophily, or other stable emergent laws & \textsf{REAL} / \textsf{INCOMPLETE} \\
Counterfactual response & Replicated randomized arms; verified intervention path; placebo and negative controls; uncertainty & Four paired event/baseline trees (v1) and the parallel-worlds engine (v2) demonstrate comparison mechanics & Causal policy effects, sign stability, transportability, or welfare recommendations & \textsf{REAL} / \textsf{DIAGNOSTIC} \\
Reproducibility & Immutable code/config/data/model snapshots; deterministic components; repeated reruns and failure accounting & Hashed result files, seeds/configs, per-arm completion, world and node identifiers, the diagnostic harness, and the Track A scorecard & Bitwise replay of hosted LLM calls or complete end-to-end reproduction of every archive & \textsf{REAL} / \textsf{FIXTURE} \\
Participant evaluation & Fixed scenarios; deterministic scoring; pre-registered sample size; per-game archive & Fourteen games in the Playground, the Track B scorecard, the Arena & That a playground game approximates a real social outcome or that a benchmark score is portable across cities & \textsf{REAL} / \textsf{INCOMPLETE} \\
\bottomrule
\end{tabular}
\end{table}

\subsection{Individual consistency}

v2 implements the prerequisites for a consistency study: named persistent
agents, structured state, episodic files, recall, reflection, longitudinal
logs, and a personality plugin that supplies Big-Five proxies for the
interruption probability.  These mechanisms allow within-agent questions
such as whether a stated preference survives intervening events or whether
an interview answer is traceable to stored experience.  The current memory
case does not answer those questions comparatively because three of four
arms lack phase 2 and the only complete arm contains one agent.  Generated
behavioral coherence is also not equivalent to correspondence with a real
person.  A stronger evaluation would repeat standardized probes before and
after controlled memory perturbations, separate temporal persistence from
prompt sensitivity, and compare against held-out observations where consent
and provenance permit.

\subsection{Macro correspondence}

Aggregating simulated variables is an implemented capability; calling the
aggregate realistic requires a separate calibration argument.  The audited
economy trace has five agents and one seed, while the benchmark scorecard
is a synthetic diagnostic fixture.  Neither represents a target population
or supports out-of-sample prediction.  v2 closes a v1 gap: the price level
now actually multiplies real expenditure, and the endogenous inflation rule
is a documented driver of the price level.  Future macro validation should
document the sampling frame, map each simulated construct to an external
measurement, reserve targets for testing rather than calibration, and
report uncertainty across seeds and plausible parameterizations.  Agreement
on a few marginal means would remain insufficient if joint distributions,
temporal patterns, or subgroup errors failed.

\subsection{Emergence and mechanism}

An emergent claim requires a pattern that is not merely written into a
prompt or initial condition and that recurs under controlled variation.
The present archive does not meet that standard: the network study records
no interactions, the emotion study has no state files, and the
polarization arms have unequal record counts.  The existence of
social-network, exposure, affect, indoor-room, organization, and
governance fields therefore establishes instrumentation, not a social law.
Appropriate tests would vary population size and topology, ablate candidate
mechanisms, quantify between-seed variability, and trace a macro pattern
back to micro-level events.  Where a field remains identically zero, the
first question is whether the mechanism and detector were active, not why
society produced no effect.

\subsection{Counterfactual validity}

The matched-run interface and the parallel-worlds engine are useful
scaffolds for counterfactual experiments, but the archived policy
differences are still single contrasts produced by an LLM-mediated
dynamical system.  A common seed does not guarantee common latent
randomness after prompts diverge, and a policy description does not
guarantee that an economic or mobility mechanism changed.  v2's
deterministic counterfactual layer (paired ATE, sign-flip $p$, BH $q$,
pre-event balance, placebo noise floor, heterogeneity, dose--response,
replication, and a non-overturning model interpretation) is an
improvement over v1's matched-run reducer, but it does not by itself
turn a single LLM-mediated trajectory into a real-world effect.  Before
interpreting a difference, researchers should verify the intervention's
path from event to structural state, replicate both arms over seeds,
pre-register target metrics, and include placebo events or
negative-control outcomes \citep{lipsitch2010negative}.  External policy
conclusions additionally require a defensible mapping to the population
and institutions of interest.  GAWorld v2 currently supports constructing
and auditing such a study, not replacing it.

\subsection{Reproducibility as a non-compensatory gate}

The artifact hashes used in this paper bind each numerical statement to
an audited byte sequence.  v2 extends the v1 ledger with three fields:
\texttt{world\_id}, \texttt{node\_id}, and \texttt{per\_arm\_status}.
These are the minimum additional information needed to reconstruct a
run under the v2 deployment model.  They do not bind the archive to a
clean released runtime: the audit began from a dirty working tree, and
hosted providers may change behavior independently of code, seed, or
sampling parameters.  Missing state files and interrupted arms further
show that successful completion must be recorded per treatment.
Consequently, a result that cannot be reconstructed from code,
configuration, inputs, model identity, world identifier, node
identifier, per-arm completion status, and complete outputs should not
be promoted because it looks plausible or matches an external statistic.

For future releases, the minimum reproducibility unit should contain the
exact code revision plus a patch or clean tag, environment lockfile,
initial profiles and maps, provider/model identifiers, prompt and
sampling settings, random seeds, per-arm status, logs, raw and derived
artifacts, the world identifier, the node identifier (if any), the
player intervention trail (if any), and a hash manifest.  Hosted-model
studies should estimate replay variability rather than promise bitwise
determinism.  The diagnostic benchmark can then serve as a regression
gate: fixtures test parsers and scoring logic, while separately labeled
simulation runs test whether scientific evidence has actually been
produced.

\subsection{Participant evaluation}

The Agent Playground and the Agent Arena produce a third class of
evidence: per-game or per-question response archives from named
residents.  The Track B scorecard combines several games under
pre-registered laws.  Participant evaluation is a distinct claim level
because it concerns whether the residents' responses under a fixed
scenario follow a rule, not whether the residents approximate a real
population.  v2 therefore requires that a Track B law gate on a
pre-registered minimum sample size; laws that fail the gate are
reported as \emph{withheld}, not as \emph{reversed} or
\emph{supported}.  The Arena is a per-resident benchmark and does not
produce a population claim; its archives are reusable but the
researcher is responsible for the population inference.

\subsection{Intended use}

At its present evidence level, GAWorld is best used for mechanism
prototyping, trace inspection, hypothesis generation, controlled software
experiments, the design of more rigorous artificial-society studies, and
the teaching of agent-based social simulation.  The archive does not
support deploying GAWorld as a substitute population, forecasting real
social outcomes, ranking policies, or making decisions about
identifiable people.  Those are not merely stronger wordings of the
same result; they are different claims requiring external data,
identification, uncertainty quantification, and governance that the
current artifacts do not provide.
